\scanpage{199}
\mspage{115}
% archive sheet number: 198
% element: text | top, line 1 | où
où
% element: display | top, lines 2-3 | pi0 factorial
\[
\begin{aligned}
\pi_0!&=\operatorname{card}\mathrm{Aut}\unsure{E_0)|(E_0)},\qquad E_0\text{ ens de cardinal }\pi_0\\
&=\pi_0^{\pi_0}=2^{\pi_0}
\end{aligned}
\]
% element: text | upper third, line 4 | donc
donc
% element: display | upper third, line 4 | card J' times two
\[
\operatorname{card}J'\cdot2^{\pi_0}=\pi^{\pi_0},
\]
% element: text | upper third, line 5 | or
or
% element: deletion | upper third, lines 6-7 | ce qui implique
\deleted{ce qui \unsure{implique}, \deleted{\unsure{pour}} si $\pi^{\deleted{\unsure{\pi_0}}}>2^{\pi_0}$, d'où $\pi^{\pi_0}>2^{\pi_0}$, que l'on a $\operatorname{card}J'>2^{\pi_0}$, \unsure{donc}}
% element: text | middle, lines 8-13 | D'autre part
D'autre part il est immédiat que $\operatorname{card}J'\geq\operatorname{card}\unsure{\mathscr P|\mathfrak P|\mathrm P}(E_0)=2^{\pi_0}$ (si $E_0$ est une partie de $S$ de cardinal $\leq\pi_0$), d'où le \deleted{\unsure{deuxième}} premier mb de ces égalités $\unsure{=|\text{à}}\operatorname{card}J'$, donc on trouve
% element: display | middle, line 14 | card J'
\[
\operatorname{card}J'=\pi^{\pi_0},
\]
% element: margin note | left beneath calculation | sorry!
\margin{\unsure{sorry!}}
% element: drawing | centre beneath calculation | short oblique separating stroke
\figure{\begin{tikzpicture}\draw (0,0.2)--(0.35,0);\end{tikzpicture}}

% element: text | lower half, lines 15-21 | Finalement
Finalement, je me suis convaincu que la prop. \unsure{4.5.3|4.5.2}, et les Corollaires et lemmes précédents par lesquels j'ai essayé de l'établir, \unsure{n'est vrai} que pour \deleted{\unsure{n'admet que pour les cardinaux}} \inserted{les cardinaux $\pi>\pi_0$ tels que} \deleted{$\pi\unsure{\geq|>}\pi_0$}
% element: display | lower half, line 21 | pi power pi0 equals pi
\[
\pi^{\pi_0}=\pi,
\]
% element: text | foot, lines 22-23 | ce qui moyennant
ce qui, moyennant l'hypothèse des continus \unsure{généralisée|généralisés}, équivaut à :
% element: display | foot, line 24 | pi power pi0 relation to two power pi
\[
\pi^{\pi_0}\unsure{<|\neq}2^\pi\ (=\pi^\pi)
\]
% element: text | foot, last line | ou mieux
(ou mieux : $\pi^{\pi_0}\neq2^\pi\ (=\pi^\pi)$). Je vais donner
